\section{Extended Theoretical Development}
\label{sec:extended-theory}

\subsection{Exact value functions and recovery maps}
For one component, let $c=\mu A^\top x$ and let $S_\gamma(c)$ denote entrywise soft thresholding. The following theorem fixes both the value function and the optimizer, including the inactive case.

\begin{theorem}[Loading-ball elimination]
\label{thm:a-loading-elimination}
For every $c\in\R^n$ and $\gamma\ge0$,
\[
\max_{\norm{z}_2\le1}\{c^\top z-\gamma\norm{z}_1\}=\norm{S_\gamma(c)}_2.
\]
If $S_\gamma(c)\neq0$, the unique maximizer is
\[
z^\star=\frac{S_\gamma(c)}{\norm{S_\gamma(c)}_2}.
\]
If $S_\gamma(c)=0$, then $z^\star=0$ is optimal; it is the unique optimum whenever $\abs{c_i}<\gamma$ for every $i$.
\end{theorem}
\begin{proof}
For any feasible $z$,
\[
c^\top z-\gamma\norm{z}_1
=\sum_i (c_i z_i-\gamma\abs{z_i})
\le \sum_i (\abs{c_i}-\gamma)_+\abs{z_i}.
\]
Choosing the signs of $z_i$ to agree with $c_i$ makes the inequality tight. Cauchy--Schwarz then gives
\[
\sum_i (\abs{c_i}-\gamma)_+\abs{z_i}
\le \norm{S_\gamma(c)}_2\norm{z}_2
\le \norm{S_\gamma(c)}_2.
\]
When the thresholded vector is nonzero, equality in Cauchy--Schwarz requires the stated normalized vector. If the thresholded vector vanishes, the objective is nonpositive and $z=0$ attains zero. Strict inequalities $\abs{c_i}<\gamma$ force every nonzero coordinate to make a strictly negative contribution.
\end{proof}

\begin{theorem}[Quadratic-lift representation]
\label{thm:a-quadratic-lift}
For every $c\in\R^n$,
\[
\norm{S_\gamma(c)}_2^2
=\max_{z\in\R^n}\left\{2c^\top z-\norm{z}_2^2-2\gamma\norm{z}_1\right\},
\]
and the unique maximizer is $z^\star=S_\gamma(c)$.
\end{theorem}
\begin{proof}
The objective is separable. For coordinate $i$, maximize
$2c_i z_i-z_i^2-2\gamma\abs{z_i}$. Restricting to the sign of $c_i$ and completing the square gives the maximizer $\operatorname{sign}(c_i)(\abs{c_i}-\gamma)_+$ and optimal value $(\abs{c_i}-\gamma)_+^2$. Summation proves the claim.
\end{proof}

\begin{proposition}[The two block aggregations are not equivalent]
\label{prop:a-non-equivalence}
For $m>1$, maximizing $\sum_j r_j(X)$ and maximizing $\sum_j r_j(X)^2$ need not have the same maximizers, even when every $r_j$ is nonnegative and arises from a thresholded norm.
\end{proposition}
\begin{proof}
It suffices to construct two feasible points with response vectors $(r_1,r_2)=(2,0)$ and $(\sqrt{3/2},\sqrt{3/2})$. The first has larger squared sum, $4>3$, while the second has larger linear sum, $2\sqrt{3/2}>2$. Such response pairs are realized by a diagonal two-dimensional data operator and two orthogonal score frames with different rotations. A concrete numerical instance is included in Appendix~\ref{app:a-detailed} and the released regression test.
\end{proof}

\subsection{Differentiability and active-region smoothness}
The scalar function $q_\gamma(u)=(\abs{u}-\gamma)_+^2$ is continuously differentiable, with
$q_\gamma'(u)=2\operatorname{sign}(u)(\abs{u}-\gamma)_+$. Hence
\[
\nabla_{x_j}F_{\mathrm{sq}}(X)=2\mu_j A S_{\gamma_j}(\mu_j A^\top x_j).
\]
It is globally Lipschitz in the ambient space with constant at most $2\mu_j^2\norm{A}_2^2$ per column because soft thresholding is nonexpansive. On any fixed support-and-sign cell, the gradient is affine and the Hessian is constant. This piecewise-quadratic structure is the basis for finite identification and post-identification linearization.

\begin{lemma}[Global Lipschitz gradient]
\label{lem:a-lipschitz-gradient}
For $f_j(x)=\norm{S_{\gamma_j}(\mu_j A^\top x)}_2^2$,
\[
\norm{\nabla f_j(x)-\nabla f_j(y)}_2
\le 2\mu_j^2\norm{A}_2^2\norm{x-y}_2.
\]
\end{lemma}
\begin{proof}
Use the gradient formula, submultiplicativity, and nonexpansiveness of soft thresholding:
\[
\norm{A[S_\gamma(\mu A^\top x)-S_\gamma(\mu A^\top y)]}_2
\le \norm{A}_2\mu\norm{A^\top(x-y)}_2.
\]
Multiplying by the leading factor $2\mu$ proves the result.
\end{proof}

\subsection{Safeguarded ascent and whole-sequence convergence}
The historical polar step maximizes the ambient affine minorant. The analyzed algorithm (A-SAFE) uses the Riemannian gradient $\eta_k=\mathrm{grad}_R F(X_k)=P_{T_{X_k}\!\St}(\nabla F(X_k))$, an explicit retraction $R_{X_k}$ satisfying standard second-order smoothness, and an Armijo acceptance rule
\begin{equation}
F\bigl(R_{X_k}(t\eta_k)\bigr)\ge F(X_k)+\sigma\,t\,\norm{\eta_k}_F^2,\qquad \sigma\in(0,1),
\label{eq:a-armijo}
\end{equation}
with geometric backtracking $t\gets\rho t$, $\rho\in(0,1)$. The next two lemmas are the paper-specific ``sufficient-ascent'' and ``relative-error'' inequalities on which the whole-sequence theorem is built.

\begin{lemma}[A6, sufficient ascent]
\label{lem:a-sufficient-ascent}
Let $F$ have $L$-Lipschitz ambient gradient (Lemma~\ref{lem:a-lipschitz-gradient}) and let the retraction $R$ satisfy, uniformly on compact subsets of $\St(p,m)$: $R_X(0)=X$, $DR_X(0)=\mathrm{Id}_{T_X\!\St}$, and $\norm{R_X(\eta)-X-\eta}_F\le\tfrac{C_R}{2}\norm{\eta}_F^2$ for $\norm{\eta}_F\le\bar t$. Then there exists $\underline t>0$, independent of $k$, such that Armijo backtracking from $t=\bar t$ selects a step $t_k\ge\underline t$, and the accepted iterate $X_{k+1}=R_{X_k}(t_k\eta_k)$ satisfies
\begin{equation}
F(X_{k+1})-F(X_k)\ \ge\ c_1\,\norm{X_{k+1}-X_k}_F^2,
\qquad c_1=\frac{\sigma\,\underline t}{(1+\tfrac{C_R}{2}\bar t\,M)^{2}\bar t}>0,
\label{eq:a-sufficient-ascent}
\end{equation}
where $M=\sup_k\norm{\eta_k}_F$ is finite because the iterates stay in a compact set.
\end{lemma}
\begin{proof}
By the retraction second-order estimate and the $L$-Lipschitz-gradient descent lemma applied to $-F$,
\[
F\bigl(R_{X_k}(t\eta_k)\bigr)
\ge F(X_k)+t\norm{\eta_k}_F^2-\tfrac{L}{2}\norm{R_{X_k}(t\eta_k)-X_k}_F^2
\ge F(X_k)+\Bigl(t-\tfrac{L\,C_R^2}{2}t^2\Bigr)\norm{\eta_k}_F^2.
\]
The Armijo bound \eqref{eq:a-armijo} is therefore satisfied whenever
$t-\tfrac{L\,C_R^2}{2}t^2\ge\sigma\,t$, i.e.\ $t\le\tfrac{2(1-\sigma)}{L\,C_R^2}$. Hence backtracking terminates at some $t_k\ge\underline t:=\min\bigl\{\bar t,\ \rho\cdot\tfrac{2(1-\sigma)}{L\,C_R^2}\bigr\}>0$. The step norm obeys $\norm{X_{k+1}-X_k}_F\le \bigl(1+\tfrac{C_R}{2}t_k\norm{\eta_k}_F\bigr)\,t_k\norm{\eta_k}_F\le (1+\tfrac{C_R}{2}\bar t M)\,t_k\norm{\eta_k}_F$. Combining with \eqref{eq:a-armijo} at the accepted step and bounding $t_k\le\bar t$,
\[
F(X_{k+1})-F(X_k)\ \ge\ \sigma\,t_k\,\norm{\eta_k}_F^2
\ \ge\ \frac{\sigma\,\underline t}{(1+\tfrac{C_R}{2}\bar t M)^{2}\,\bar t}\,\norm{X_{k+1}-X_k}_F^2.
\]
\end{proof}

\begin{lemma}[A7, relative error]
\label{lem:a-relative-error}
Under the same hypotheses, the limiting subdifferential of $\Psi:=-F+\iota_{\St}$ at $X_{k+1}$ contains an element
\[
w_{k+1}\in\partial\Psi(X_{k+1})
\quad\text{with}\quad
\norm{w_{k+1}}_F\ \le\ c_2\,\norm{X_{k+1}-X_k}_F,
\]
where $c_2$ depends only on $L$, $C_R$, $\sigma$, and the compact set containing the iterates. An admissible choice is $w_{k+1}=-\mathrm{grad}_R F(X_{k+1})=-\nabla F(X_{k+1})+\mathrm{proj}_{N_{\St}(X_{k+1})}\nabla F(X_{k+1})\in\partial(-F)(X_{k+1})+N_{\St}(X_{k+1})\subseteq\partial\Psi(X_{k+1})$.
\end{lemma}
\begin{proof}
Since $F$ is smooth, $\partial(-F)(X)=\{-\nabla F(X)\}$; the Riemannian gradient equals the ambient gradient minus its projection onto the normal cone of $\St(p,m)$ at $X$, so $-\mathrm{grad}_R F(X_{k+1})\in\partial\Psi(X_{k+1})$. Triangle inequality:
\[
\norm{\mathrm{grad}_R F(X_{k+1})}_F
\le \norm{\mathrm{grad}_R F(X_{k+1})-\mathrm{grad}_R F(X_k)}_F+\norm{\eta_k}_F.
\]
On a compact subset of $\St(p,m)$, the Riemannian-gradient field inherits a Lipschitz constant $L'=L\cdot(1+\kappa)$ where $\kappa$ bounds the derivative of the tangent-projection map; the first term is at most $L'\norm{X_{k+1}-X_k}_F$. For the second, Lemma~\ref{lem:a-sufficient-ascent} gives $\norm{\eta_k}_F\le\norm{X_{k+1}-X_k}_F/\underline t$. Hence
\[
\norm{w_{k+1}}_F\ \le\ \Bigl(L'+\tfrac1{\underline t}\Bigr)\,\norm{X_{k+1}-X_k}_F\ =:\ c_2\,\norm{X_{k+1}-X_k}_F.
\]
\end{proof}

\begin{theorem}[A8, whole-sequence convergence under K\L]
\label{thm:a-whole-sequence}
Let $(X_k)_{k\ge0}\subset\St(p,m)$ be generated by the safeguarded Riemannian ascent of Section~\ref{sec:method}, with retraction satisfying the second-order estimate of Lemma~\ref{lem:a-sufficient-ascent}. Then
\begin{enumerate}
\item\label{item:a-whole-finite-length} the sequence has finite length, $\sum_{k\ge0}\norm{X_{k+1}-X_k}_F<\infty$;
\item\label{item:a-whole-cauchy} $(X_k)$ converges to a single Riemannian critical point $X_\star$ of $F$ on $\St(p,m)$;
\item\label{item:a-whole-rate} with K\L{} exponent $\theta\in[0,1)$ of $\Psi=-F+\iota_{\St}$ at $X_\star$, the rate is finite if $\theta=0$, linear $\norm{X_k-X_\star}_F\le C\varrho^k$ (some $\varrho\in(0,1)$) if $\theta\in(0,\tfrac12]$, and $\norm{X_k-X_\star}_F\le C k^{-(1-\theta)/(2\theta-1)}$ if $\theta\in(\tfrac12,1)$.
\end{enumerate}
\end{theorem}
\begin{proof}
Compactness of $\St(p,m)$ makes the iterates bounded and $F$ bounded above. Lemma~\ref{lem:a-sufficient-ascent} rewrites as the sufficient-decrease inequality $\Psi(X_k)-\Psi(X_{k+1})\ge c_1\norm{X_{k+1}-X_k}_F^2$ for $\Psi=-F+\iota_{\St}$; Lemma~\ref{lem:a-relative-error} supplies the relative-error inequality in the form required by \citet{attouch2013convergence}.

Soft-thresholding is piecewise affine on a finite polyhedral partition, hence semi-algebraic; a sum of squared Euclidean norms of composite semi-algebraic maps is semi-algebraic, so $F$ and $-F$ are semi-algebraic. The Stiefel manifold $\{X:X^\top X=I_m\}$ is a real algebraic variety, so $\iota_{\St}$ is semi-algebraic. Sums of semi-algebraic functions are semi-algebraic, so $\Psi$ is semi-algebraic and therefore definable in the o-minimal structure of real semi-algebraic sets. Definable functions satisfy the Kurdyka--\L{}ojasiewicz inequality on their domain \citep{bolte2007lojasiewicz}; in particular at every limit point of $(X_k)$.

The hypotheses of the abstract convergence theorem of \citet{attouch2013convergence} (sufficient decrease, relative error, continuity along the sequence, boundedness of $(X_k)$, and K\L) are therefore met for $\Psi$. Applying it yields items~\ref{item:a-whole-finite-length} and~\ref{item:a-whole-cauchy}. The rate claim \ref{item:a-whole-rate} is the standard K\L-exponent corollary of the same framework.
\end{proof}

The constants $c_1,c_2$ in Lemmas~\ref{lem:a-sufficient-ascent}--\ref{lem:a-relative-error} are computable from the Lipschitz bound (Lemma~\ref{lem:a-lipschitz-gradient}) together with the retraction curvature $C_R$ and the Armijo parameters; the reference implementation records them in the solver trace. The rank-deficient polar branch is handled by the deterministic SVD completion of Section~\ref{sec:setup} and does not alter the above constants on the compact iterate set.

\subsection{Identification with explicit margins}
Let $\mathcal I_j(X)=\{i:\abs{\mu_j a_i^\top x_j}>\gamma_j\}$ and let $\mathcal S_j(X)$ denote active signs. The strict margin is
\[
\delta_\star=\min_{i,j}\left|\abs{\mu_j a_i^\top x_{j\star}}-\gamma_j\right|.
\]
The conditional identification theorem in Section~\ref{sec:theory} is immediate from continuity. A stronger perturbation statement is useful for screening.

\begin{corollary}[Support stability radius]
\label{cor:a-support-radius}
If $\delta_\star>0$ and
\[
\norm{x_j-x_{j\star}}_2<\min_i\frac{\delta_\star}{2\mu_j\norm{a_i}_2},
\]
then component $j$ has the same active support as at $X_\star$. If active correlations are also separated from zero by $\sigma_\star>0$, replacing the numerator by $\min(\delta_\star,\sigma_\star)/2$ preserves active signs.
\end{corollary}
\begin{proof}
Cauchy--Schwarz bounds every correlation perturbation by $\mu_j\norm{a_i}_2\norm{x_j-x_{j\star}}_2$. No threshold or sign boundary can be crossed inside the stated radius.
\end{proof}

\subsection{Common-support block rate}
When all columns share an active set $S$ in the hard-threshold model, the fixed-support update is orthogonal iteration for $C_S=A_SA_S^\top$.

\begin{theorem}[Common-support subspace contraction]
\label{thm:a-block-rate}
Let $C_S$ be symmetric positive semidefinite with $\lambda_m(C_S)>\lambda_{m+1}(C_S)$. Let $U_\star$ span the dominant $m$-dimensional invariant subspace, and assume $U_\star^\top X_0$ is nonsingular. Orthogonal iteration satisfies
\[
\norm{\tan\Theta(X_{k},U_\star)}_2
\le
\left(\frac{\lambda_{m+1}(C_S)}{\lambda_m(C_S)}\right)^k
\norm{\tan\Theta(X_{0},U_\star)}_2.
\]
\end{theorem}
\begin{proof}
Write the initial frame in the eigenbasis of $C_S$ and partition it into dominant and complementary blocks. After multiplication by $C_S^k$, the complementary-to-dominant coefficient ratio is multiplied by at most $(\lambda_{m+1}/\lambda_m)^k$. Orthonormalization does not change the represented subspace, yielding the tangent-of-principal-angle bound.
\end{proof}

\subsection{Heterogeneous supports and the polar Jacobian}
For component-specific supports, define $G(X)=[C_{S_1}x_1,\ldots,C_{S_m}x_m]$ and $\Phi(X)=\operatorname{Polar}(G(X))$. If $G(X_\star)$ has full column rank, the polar map is differentiable. Its differential is characterized by a Sylvester equation for the tangent skew term. The correct local condition is therefore
\[
\rho\left(D\Phi(X_\star)|_{T_{X_\star}\St}\right)<1,
\]
not one pooled eigengap. Appendix~\ref{app:a-polar} derives the operator needed to evaluate this condition numerically.

\subsection{Second-order statements that are actually justified}
On a fixed rank-one support, the objective reduces to a Rayleigh quotient. A simple nonleading eigenvector has a tangent direction of positive curvature toward a larger eigenvalue and is not a local maximum. For heterogeneous supports and threshold boundaries, no generic strict-saddle claim is made; this paper only reports second-order classification for fixed-support smooth cells whose tangent Hessian is explicitly computable.


\begin{figure}[t]
\centering
\begin{tikzpicture}[node distance=7mm and 9mm]
\node[primarybox,text width=2.55cm] (n1) {Exact model};
\node[primarybox,text width=2.55cm,right=of n1] (n2) {Specified algorithm};
\node[secondarybox,text width=2.55cm,right=of n2] (n3) {Global convergence};
\node[secondarybox,text width=2.55cm,below=of n3] (n4) {Support identification};
\node[secondarybox,text width=2.55cm,left=of n4] (n5) {Local dynamics};
\node[successbox,text width=2.55cm,left=of n5] (n6) {Controlled + neural evidence};
\draw[arrPrimary] (n1)--(n2);
\draw[arrPrimary] (n2)--(n3);
\draw[arrPrimary] (n3)--(n4);
\draw[arrPrimary] (n4)--(n5);
\draw[arrPrimary] (n5)--(n6);
\end{tikzpicture}
\caption{Theory-to-evidence chain for the developed manuscript.}
\label{fig:a-extended-chain}
\end{figure}

\begin{table}[t]
\centering
\caption{Extended theorem, proof, and evidence obligations.}
\label{tab:a-extended-obligations}
\small
\begin{tabular}{p{0.20\linewidth}p{0.30\linewidth}p{0.39\linewidth}}
\toprule
Layer & Required statement & Acceptance evidence \\
\midrule
Algebra & Exact value functions and recovery maps & Symbolic derivation plus numerical identity tests \\
Global dynamics & Safeguarded whole-sequence result & Sufficient-ascent and relative-error lemmas for implemented retraction \\
Active structure & Finite support/sign identification & Strict-margin theorem and tie counterexample \\
Local theory & Common- and heterogeneous-support rates & Eigengap proof or polar-Jacobian contraction diagnostic \\
Application & LLM activation subspaces & Split stability and causal interventions \\
\bottomrule
\end{tabular}
\end{table}

